
\statement{C.39}{Proved}{Every peripheral involution passes the symmetric-four test and the dyadic relative \(S_3\) levels}
Let \(\alpha=(-1,f)\) be an involution in the full peripheral
group \(\mathcal P\), with its normalized representative and
full R1 word \(u\in F'\), so that \(u=f\alpha(u)\).
For every epimorphism \(q:F\to S_4\) and every
\(\mu\in\widehat{\mathbb Z}^{\times}\),
\[
        \bigl\langle q(x^\mu),q(u^{-1}y^\mu u)\bigr\rangle=S_4.
 \tag{C.39.1}
\]
For every epimorphism \(q_3:F\to S_3\), put \(U=\ker q_3\).
The normalized substitution
\[
       \theta_\mu:x\longmapsto x^\mu,\qquad
                 y\longmapsto u^{-1}y^\mu u
\]
induces an automorphism of the actual relative pro-two group
\[
                              F/O^2(U),
 \tag{C.39.2}
\]
and intertwines the induced \(\alpha\) and \(c\).
Neither duality nor a hexagon or pentagon is assumed. These are
generation and unmarked conjugacy statements; relative third-cusp
preservation is not asserted.

\begin{proof}
\emph{The symmetric-three head.}\par
Write \(\psi=\theta_1\) and \(y'=u^{-1}yu\).
For an epimorphism to \(S_3\), the images of \(x,y\)
are either two distinct reflections, or a nonidentity rotation
and a reflection in one of the two orders. The reflection case
is the \(p=3\) case of A.44. Its proof uses only the full
peripheral involution, the full R1 equation, and the section
incidence of C.28; it uses no \(\widehat{\mathrm{GT}}_0\)
symmetry relation. In the mixed cases \(q_3(u)\) is in the
rotation group \(S_3'=C_3\). If \(y\) is rotational its
image is unchanged by this conjugation; if \(y\) is a
reflection, it is still a reflection while the nonidentity
rotational image of \(x\) is retained. Thus
\[
                              q_3\psi(F)=S_3
 \tag{C.39.3}
\]
in every case.

Each of the three ordered generator-type possibilities is a
single \(\operatorname{Aut}(S_3)\)-orbit of generating pairs.
Therefore \(q_3\psi\) has the same kernel as \(q_3\).
The same conclusion holds for \(q_3\theta_\mu\), since unit
powers preserve orders and generated cyclic subgroups, and
for \(q_3\alpha,q_3c\). In particular \(U\) is preserved
by all the substitutions and involutions under consideration.
Its pro-two residual is fully invariant under endomorphisms
of \(U\), so these maps induce endomorphisms of (C.39.2).

\emph{The symmetric-four test.}\par
Let \(q:F\to S_4\) be onto, and write
\(a=q(x)\), \(b=q(y)\), \(b'=q(y')\).
The quotient by the normal Klein group gives
\(S_4/V_4\simeq S_3\); hence (C.39.3) shows that
\(M=\langle a,b'\rangle\) still maps onto \(S_3\).

Every proper subgroup \(M<S_4\) mapping onto \(S_3\)
is a point stabilizer \(S_3\). Indeed \(MV_4=S_4\), and
\(M\cap V_4\) is stable under \(M\) and under the
abelian group \(V_4\), hence under all of \(S_4\).
The natural \(S_3\)-module \(V_4\) is irreducible, so
this intersection is either \(1\) or \(V_4\).
The latter would give \(M=S_4\); thus \(|M|=6\).
A subgroup of order six in \(S_4\) fixes a point: its
orbits cannot have size four, and a faithful action with
only orbits of size at most two would have abelian image.
It is the full stabilizer of that point and is self-normalizing.

If either original generator has order four, the substituted
pair still contains an element of order four, and cannot lie
in such a point stabilizer. Otherwise a generating pair of
\(S_4\) has orders two and three in some order: two
involutions generate a dihedral group of order at most eight
in \(S_4\), and two three-cycles generate an even subgroup.
The order-two element is a transposition, since a double
transposition and a three-cycle would both be even. A
transposition and a three-cycle generate \(S_4\) exactly
when their supports have union all four letters. Their
product is then a four-cycle. These pairs form one inner
automorphism orbit of passport \((2,3,4)\), and their
interchanges form one orbit of passport \((3,2,4)\).

For either critical orbit, \(q\alpha\) is another generating
pair of the same ordered types, since \(\alpha\) is a full
automorphism. It therefore has the same kernel as \(q\);
the induced action is \(\operatorname{Inn}(h)\) for some
\(h\in S_4\). Since \(\alpha^2=1\) and \(S_4\) is
center-free, \(h^2=1\). The canonical equation gives
\(\alpha(x)=x^{-1}\) and \(\alpha(y')=y'^{-1}\);
thus \(h\) inverts \(a,b'\).

If \(M\) were proper, its self-normalization would put
\(h\) in this \(S_3\) point stabilizer. There it
centralizes the order-two generator and inverts the
order-three generator. Consequently it equals the
order-two generator. In the finite extension map
\[
 F\rtimes_\alpha C_2\longrightarrow
 S_4\rtimes_{\operatorname{Inn}(h)}C_2
                   \simeq S_4\times C_2
\]
one of the full sections \(x\alpha,y'\alpha\) therefore
has central image \((1,c)\). The third cusp joins these
two section classes in the actual incidence triangle
of C.28. Its two local reflections have images whose
product has order at most two, because one is central.
The actual image of the third cusp has order four in
both critical passports. This contradiction proves
(C.39.1) for \(\mu=1\), and unit powering preserves the
generated subgroup, giving the assertion for every \(\mu\).

The exact finite facts are separately verified in
\path{work_CE/code/check_C39_S4.py}: all thirty subgroups,
all \(216\) generating pairs in nine automorphism orbits,
and the two critical orbit/normalizer/inversion checks.
After imposing the already proved \(S_3\)-generation
condition, exactly twelve proper-image candidates remain,
all with the central-section contradiction just used.

\emph{The actual relative pro-two group.}\par
Fix \(U=\ker q_3\), and put \(G=F/O^2(U)\) and
\(N=U/O^2(U)\). Thus \(N\) is free pro-two of rank
seven and \(G/N=S_3\). We have already proved that
\(\theta_\mu\) induces an endomorphism of \(G\) and
an automorphism of this finite head.

Suppose its closed image is proper. There is then a
finite quotient \(\pi:G\to D\), retaining the head
\(S_3\), for which
\[
 H=\pi\theta_\mu(G)<D,\qquad P=\pi(N)\text{ is a two-group},
 \qquad HP=D,\qquad D/P=S_3.
\]
No endomorphism of this particular \(D\) is assumed.
The subgroup \(H\Phi(P)\) remains proper. Otherwise
\(P=(H\cap P)\Phi(P)\), and the finite two-group
Frattini property would give \(H\cap P=P\), hence
\(H=D\).

Pass to \(\overline D=D/\Phi(P)\).
Its kernel \(\overline P\) is an elementary abelian
two-group. The subgroup
\(K=\overline H\cap\overline P\) is normal in
\(\overline D=\overline P\,\overline H\):
\(\overline H\) normalizes its intersection and the
abelian \(\overline P\) centralizes it.
Consequently \(\overline D/K\) is a genuinely split
extension of the nonzero \(\mathbb F_2[S_3]\)-module
\(\overline P/K\), with complement
\(\overline H/K\simeq S_3\). Choose a nonzero simple
module quotient \(W\). We obtain an actual quotient
\[
                                  W\rtimes S_3
 \tag{C.39.4}
\]
of \(G\), in which the image of \(\theta_\mu\) is
contained in a complement.

There are precisely two simple
\(\mathbb F_2[S_3]\)-modules: the trivial one-dimensional
module and the natural two-dimensional module. This
follows by restricting to the normal odd subgroup \(C_3\).
Its trivial constituent gives a simple module for \(C_2\)
in characteristic two, necessarily trivial; its unique
nontrivial irreducible constituent over \(\mathbb F_2\)
is the two-dimensional one, on which the quotient \(C_2\)
interchanges the two cubic characters.

For the trivial module, (C.39.4) is \(C_2\times S_3\).
Its abelianization has mod-two rank two, while a complement
has abelian image of rank one. This is impossible because
\(\theta_\mu\) induces identity on \(H_1(F,\mathbb F_2)\):
\(u\in F'\) and every unit \(\mu\) is odd.
For the nontrivial module, (C.39.4) is
\(\mathbb F_2^2\rtimes S_3\simeq S_4\), contradicting
(C.39.1). These exhaust the simple quotients, so the
closed image of \(\theta_\mu\) in \(G\) is all of \(G\).
The finitely generated profinite group \(G\) is Hopfian,
and this endomorphism is an automorphism.

Finally the full equation \(u=f\alpha(u)\) gives
\(\alpha\theta_\mu=\theta_\mu c\) on \(x,y\).
All indicated kernels are preserved, so this identity
descends to (C.39.2). It proves the stated unmarked
relative conjugacy and does not assert a third-cusp identity.
\end{proof}
